%MANUSCRIPT: Description of Source Program GRAZING_9
\newcommand{\eto}{{\rm e}}                            % e, exp(1)
\newcommand{\half}{\mbox{$\frac{1}{2}$}}           % "Small" 1/2
\documentstyle[12pt]{article}
\def\today{\bf \ifcase\month\or
January\or February\or March\or April\or May\or June\or
July\or August\or September\or October\or November\or December\fi
\space\number\day, \number\year}
\textwidth 149mm
\textheight 245mm

\begin{document}
\thispagestyle{empty}
\topmargin 0pt
\headheight 0pt
\headsep 0pt
\topskip 9mm
\oddsidemargin 5mm

%\baselineskip=20pt

\begin{center}
{\large{\normalsize\sf Draft as of \today ~~File: instruction.tex}\\[3ex]}
 {\Large \bf{Program GRAZING}\\{\large version 9}\\{\normalsize A Fortran Program for 
  Estimating Reactions\\ in Collision between Heavy Nuclei.}}
\end{center}

\vskip 5truecm
\begin{center}
by
\vskip 2truecm
Aage Winther\\
The Niels Bohr Institute\\
Blegdamsvej 17\\
DK-2100 Copenhagen �, Denmark\\
\vskip 2truecm
\end{center}

\begin{abstract}
The interactive Fortran program described below evaluates the 
reactions to be expected
in a collision between two heavy nuclei at moderate bombarding energies.
It calculates the cross section for the 
distribution in mass and charge, kinetic energy and excitation energy
of the two scattered nuclei. It also calculates the cross section 
leading to capture including coupled channel effect in the tunneling
of the fusion barrier. \end{abstract}
\newpage

\section{Introduction}\label{intr}

The program GRAZING described below is created in order to evaluate the
observables in the multinucleon transfer reactions that happen in the
region outside the sum of the half-density radii of two colliding heavy
nuclei, i.e. in the grazing region. The program is not fully developed and 
the present version GRAZING\_9 certainly still contain some bugs. I would 
be happy to receive information about these in order to be able to
correct them in  future versions.
\par
The calculations follow closely the notation and formulae given in 
the two papers \cite{W1} and \cite{W2}.The default parameters from nuclear
structure are average parameters based on our theoretical and
semiempirical knowledge of nuclear radii, formfactors, level
densities, etc. and need adjustement for individual cases. It has
until now not been systematically compared to experimental data.
\par
The program is highly modular, such that the evaluation of a given
quantity or a given operation is always made in the same subroutine
or function routine. The general routines are collected in three 
libraries: The general mathematical routines, MATH\_LIB, the angular
momentum algebra, ANGULARMOM\_LIB, the general physics routines, 
FYS\_LIB, and the standard routines related to the trajectory, 
AKYUZ\_LIB. These files contain a few routines that are not used in 
the present program. If an error occurs in any subroutine the flag 
{\it ierr} is set equal to 1. Other values ($\not= 0$) may indicate 
reduced accuracy etc. The programs make extensive use the NAG library. 
\par
The structure of the program is shown in Fig.\ref{fig1}. It consist
of a main program, where the desired type of calculation is specified,
and 8 legs that perform a specified job. The main program further
provides facilities for calculations of total cross sections and 
excitation functions, and calculates quantities needed in the LEGs.
\par
An important element in the program is the determination of the
capture distance, which is mainly governed by the low-lying collective 
modes. This distance, {\it rcapt}, is calculated in a subroutine, 
CAPTURE\_S, which essentially follows ref.\cite{BDW} in solving 
the classical equation of motion for the relative position of the 
two ions coupled to the collective modes through the surface-surface 
interaction (i.e. neglecting transfer but including Coulomb excitation). 
\par
The legs are numbered in increasing complexity of the calculation.
Thus the subroutine LEG1\_S only calculates mean values of the observables
and their spread. LEG2\_S calculates the distribution in N and Z, and the 
mean values and spreads of the excitation energy and angular momentum 
for each partition. LEG3\_S calculates the probability distribution 
in energy as well as the mean angular momentum and drift in mass and 
charge for given energy loss. LEG4\_S calculates the energy distribution 
for each partition as well as the mean angular momentum loss for given
energy loss and given partition. In LEG\_4 one may also obtain the
distribution in excitation energy for specified ejectile or recoil 
and an the distribution in N and Z of ejectiles or recoils after evaporation
based on approximate excitation eneregy distributions. The distributions 
in angular momenta of the ejectile and recoil are supposed to be evaluated 
in LEG5\_S to LEG8\_S. They are not implemented.
\par
The dimensions of the arrays containing the distributions in neutrons,
protons and energy are specified in the file GRAZING.ICL. They may
be changed, if the program is afterwards recompiled.
\par
The program is constructed as an interactive program, and you are
asked all questions needed to run the program. There are three versions
of the program governed by the quantity {\it version}. \\
~1) If {\it version} = {\it vms} (one of the first statements
in the main program), the input, which you 
make will be stored in an input files, GRAZING\_IN\_xx.DAT, where you 
may also specify an identification (xx). You may then afterwards modify this 
file and run the program in an automatic mode, or use it as input in a 
batch mode. In this version output is mostly collected in named files. \\
~2) If {\it version} = {\it unix} the input which you 
make will be stored in FOR010.DAT and all output in other unnamed files.\\
~3) If {\it version} = {\it unix*} is like 1) above but works in
UNIX with named input and output files.
\par
Preparing an input file by running interactively may take a long 
time, and you may therefore instead run in a preparation mode where 
all questions are asked to prepare a correct GRAZING\_IN\_xx.DAT file, but 
since many do-loops are truncated, the numerical results are useless.
\par
The output on the screen (or in your log file) contains warning and
error messages. Most other output is found in the GRAZING\_OUT.DAT
file (in vms and unix* versions) or in FOR010.DAT in the unix version.
\par
The basic variable determining the strength of the interaction in the 
collision of the two  nuclei is the distance {\it r}(fm)of closest  
approach. It is connected to the angular momentum {\it al}
or impact parameter {\it rho}(fm) through the potential (Coulomb + nuclear)
which is assumed to be the one given in ref.\cite{AW-81}. All subroutines
directly related to the potential and the trajectory are collected
in the library AKYUZ\_LIB. If the parameters of the potential are changed 
here (in the subroutine CONSTANTS\_S) they will be changed everywhere.
A change in the radius parameter {\it vdr} is implemented as a 
standard possibility. The scattering angle {\it theta} is the elastic 
scattering angle, and does not describe the scattering angle of the 
reaction products. In the present version 9 of the program the quantity
{\it vdr} is changed during the running of the program for trajectories
with an appreciable capture probability, to take into account that the 
mean field for the scattered trajectory is different from the mean
field for elastic scattering (see below). The input value of {\it vdr} 
is stored as {\it vdr0}.
\par
Since the calculations follow refs.\cite{W1} and \cite{W2}, the
following sections mainly describe points where the program
deviates from these references because of missprints or errors, or has 
options that are not explicitly worked out. 

\section{Main Program GRAZING\_9}\label{GRAZ}
As is shown on the flowdiagram in Fig.\ref{fig1} you are asked first
(statement number 1) which type of run you intend to make (see above).
Next (sn 17) two parameters for excitation functions {\it ecmm} and
{\it decm} are specified (for no excitation functions). At this point 
you may re-enter for a new job. After specification of the colliding nuclei
the first questions asked (sn 19) is, whether the standard parameters,
given in the above references or described below, are to be used. You 
may here also change the overall accuracy of the calculations, for which
the standard value is nominally $10^{-4}$. Also one is asked about 
various options (sn 77). They include the possibitity to calculate the 
transfer reactions with an artificially narrow $Q$-value window to 
reproduce the classical picture (friction), or to calculate the final 
population of the single particle levels.
\par
The parameters for the standard potential is calculated in the subroutine 
CONSTANTS\_S (sn 76).
\par
For the inclusion of the collective modes (sn 30) one may choose to
include besides the low-lying modes, which are essential for the capture
also the giant resonances (2+ and 3-). While the properties of the former 
are taken from a table of experimental data in the subroutine QVALETC\_S
the properties of the latter are calculated semiempirically.
\par
In the subroutine QVALETC\_S (sn 30) also the Fermi surfaces of the 
colliding nuclei are determined from the separation energies for neutrons 
and protons as obtained though the MASSE\_S subroutine in the mass tables 
MASSEXP.DAT and MTABLEX.DAT containing experimental and semiempirical masses.
Various options for constructing the average $Q$-values are possible.
The $B(E2)$- and $B(E3)$-values of the low-lying collective states 
are found through the BE23C\_S subroutine in the BE23.DAT table.
\par
As is seen from Fig.\ref{fig1}(sn 31) the program after preparation of the 
basic parameters calculates the energy independent quantities, i.e. the
height {\it ebar} and radius {\it rbar} of the Coulomb barrier as well
as the critical energy {\it ecrit}, where the Coulomb barrier 
disappears, and the corresponding radius {\it rcrit}.
\par
The energy of relative motion in the center of mass system {\it ecm}
is used throughout. As input (sn 22) it may be given as the bombarding 
energy in the laboratory system {\it elab} or as the bombarding energy
per nucleon {\it epslab}. If excitation functions are to be calculated,
the initial value of {\it ecm} or {\it epslab} should be given. In this 
case the program re-enters at this point and {\it nestep} counts the number 
of steps. For higher energies a semiempirical reduction factor of the ion-ion
potential {\it zeta} is important, but the calculations are not reliable 
if $epslab>10 MeV$. The trajectory independent characteristics are 
evaluated (sn 80) in the subroutines LGRAZ\_S and RAINBOW\_S, i.e. the
grazing $l$-value {\it algraz} and the rainbow angular momentum {\it alrain}
for the reaction. 
\par
Next (sn 82) the classical capture angular momentum is determined by 
solving the coupled equations of motion for the trajectory coupled to 
the surface modes. The subroutine CAPTURE\_S includes Coulomb excitation,
but no transfer reactions. It finds capture angular momentum 
{\it alcapt} i.e. the lowest impact parameter, where the surface-surface 
distance {\it s} becomes negative, and capture distance {\it rcapt}.
As a default there is in this program included a weak non-linear 
damping of the low-lying modes. If the the energy is too low, the program 
does not find any capture, and the bombarding energy (inside the 
subroutine) is changed until capture takes place for central collisions.
The output is then the capture energy {\it ecmcapt} and {\it alcapt} = 0.
The program also calculates the orbital integrals for each magnetic
substate of the low-lying collective modes and the values of $\eta$
({\it etabs(2,3,2,4)}) determining the barrier distributions are stored
in common/capt/.
\par
The fact that the surfaces are deformed at capture means that the 
Coulomb barrier is shifted. The shift in the barrier is taken into 
account by changing the parameter {\it vdr} by an amount {\it vdrcap}
until the new barrier for {\it ecm = ecmcapt} appears at {\it al = alcapt}. 
This is achieved in the subroutine BARRIERPOS\_S (sn 69). This routine 
also evaluates the inverse barrier frequency {\it homegg} including the 
fact that the mass associated with the the barrier penetration may be 
smaller than the reduced mass of the two ions by a factor 
{\it corr = 1/(1+ delta)} due to the mass parameters of the surface 
modes.  The quantal distribution of barriers then is evaluated in
the subroutine PROBBAR\_S, which evaluates the spreading width from
transfer reactions and the distribution of barriers from the collective
modes at the capture distance. 
\par
After having chosen possible output of intermediate results (sn 71), and
determined the jobtype (sn 21), one decides whether total cross sections 
({\it itot = 1}) are to be calculated. The integration is fundamentally
performed in steps {\it delr} in {\it r}, which is the quantity that
mainly determines the probabilities for the reactions that take place.
The stepwidths are doubled after each 10 steps, and the results 
for all impact parameters are thus evaluated and integrated
until no reactions take place. For angular momenta close to 
{\it alcapt} it is , however, important to take into account that
the surfaces have been deformed and the ion-ion potential changed.
In this region it is therefore necessary to use the impact parameter
or the corresponding angular momentum {\it al} as the stepping parameter  
(in steps of {\it delal}), and get the corresponding distance of
closest approach {\it r} from the knowledge of the potential for 
the angular momentum {\it al}. The potential is calculated for given
$l$ is obtained from the subroutines TRANSM\_S and BARRIERPOS\_S.
Besides calculating the probability of capture, {\it ppcapt}, TRANSM\_S 
calculates the mean energy loss {\it descat} of those projectiles, 
that were not captured, or equivalently the mean increase in the height 
of the barrier. In BARRIERPOS\_S this is translated into a shift of the
radii {\it vdral}.
\par
If total cross sections are to be evaluated, the stepwidth in {\it r}
is calculated, and a preparation is made to determine the stepwidth in
$l$ ({\it al}) for small impact parameters. For a trajectory at the 
capture angular momentum ({\it alcapt},  there is about 50 percent 
chance that the projectile is captured. For the ones scattered the 
mean value of the barrier height is calculated with the subroutines 
TRANSM\_S and BARRIERPOS\_S , and one obtains a value for
shift {\it vdral}(see above). Evaluating the turning point {\it rtur}
for {\it al = alcapt} in the corresponding potential and the $l$-value
{\it al1} for the turning point, {\it rtur + delr} gives a stepwidth
in $l$, {\it delal = al1 - alcapt}, which is consistent with the 
stepwidth in $r$, and which is used for all impact parameters, where 
there is an appreciable probability for capture, and where the potential 
therefore depends on the impact parameter.  
\par
For total cross sections one also needs to know the dimensions for the
$N,Z$-distributions and the stepwidth in the energy-distributions to
be calculated. This information is obtained (sn 87) by running the main 
part of the program (sn 57) in a preparation mode ({\it ieprep = 1}).
In this mode one is also asked (sn 89), whether excitation functions
are to be evaluated. Since this may take long times, an upper limit
for the CPU time must be given, and a time estimate is made after the 
first energy. 
\par
The contributing trajectories are in principle specified by {\it r},
and two counters {\it istep} and {\it irstep} keep track of the number 
of steps. The integration starts with $l = 0$ and proceeds in steps
{\it delal} in $l$ for small impact parameters until the capture probability
is negligible. During this stage {\it istep} counts the number of steps
while {\it irstep} is zero. For the remaining part of the integration
(steps {\it delr} in $r$){\it irstep} counts the number of $r$-steps,   
while {\it istep} counts the total number of steps.
\par
One may instead choose to calculate a single trajectory ({\it itot = 0}) 
specified either by the angular momentum, {\it al}, or impact parameter,
{\it rho}. The trajectory can also be specified by the distance of 
closest approach for the elastic orbit or by the scattering angle for 
this trajectory.
\par
For each trajectory the program evaluates the capture probability in the 
WKB approximation in the subroutine TRANSM\_S on the basis of the 
probability distribution of the barrier calculated at the capture distance.
\par
The main program first evaluates the inverse collision time $\omega_0$
times $\hbar$ for neutron and proton transfer, {\it homegan} and 
{\it homegap}, and for inelastic scattering {\it homegain}. Since the
radial acceleration for elastic scattering close to orbiting vanishes, 
the collision time is evaluated reducing the nuclear attraction by a 
factor 0.25. This is an estimate from ref.\cite{BW} Fig.IV.19. It is 
quite uncertain, but is hopefully not very important. 
\par
In evaluating the formfactor for inelastic scattering for trajectories
reaching the classical capture distance or beyond, and being reflected,
one must realize that the low surface modes get out of control. 
The formfactors at such distances have thus been reduced with the 
argument that for those trajectories that escape capture the 
surface-surface distance has by zero-point fluctuations been 
further apart than the average estimate.
 Coulomb excitation which contributes 
very little to the energy and angular momentum loss is neglected.
With the same argument, the formfactors for transfer should be reduced.
The parameters that adjusts these reductions are uncertain, and can be 
changed in the input.
\par
The program calculates the probability for elastic scattering and the 
associated imaginary part of the optical potential. Then the 
mean-values and spreads of the six quantities  $\Delta N$, $\Delta Z$, 
$E_a^*$, $E_A^*$, $l_a$ and $l_A$ are evaluated. At this point it should
be noted that the mean excitation energy of the projectile, 
$\langle E_a^*\rangle $, and of the target, $\langle E_A^*\rangle $,
are given by
\begin{equation}
  \langle E_{a \tau}^*\rangle = \langle E_{A \tau}^*\rangle 
      = \langle E_{\tau}^*\rangle /2
  ~,  \label{eq1}
\end{equation}
while the corresponding spreads are given by
\begin{equation}
  \sigma_{a \tau}^2 = \sigma_{A \tau}^2 = \sigma_{\tau}^2 /3
  ~,  \label{eq2}
\end{equation}
where the mean values and spreads for $\tau = \nu S, \nu P,..$ 
(neutron stripping, neutron pickup etc.) are given in 
ref.\cite{W1} Eqs.(4.12) and 4.15).
\par
The program then proceeds to do the job specified in the LEG subroutines. 
The flag {\it ierr} of these subroutines indicates whether the 
integration for total cross section is finished or not, whether only 
one trajectory was to to calculated, or whether an error has occured.

\subsection{Subroutine QVALETC\_S}\label{QVAL}
In Fig.\ref{fig2} is shown the flowdiagram for the subroutine QVALETC\_S, where
most of the necessary input is prepared. It relies on the availability
of a mass tables MASSEXP.DAT and MTABLEX.DAT and a table, BE23.DAT, of 
the B(E2)- and B(E3)-values as well as excitation energies of the lowest
2+ and 3- states.
\par
The standard option for determining the Q-values for the ground state to
ground state transitions is to take the difference between the Fermi 
levels in target and projectile i.e.
\begin{equation}
   Q^{PS}_{gg} =-Q^{PP}_{gg} =  \lambda^Z_a - \lambda^Z_A \label{eq22}
\end{equation}
and
\begin{equation}
   Q^{NS}_{gg} =-Q^{NP}_{gg} =  \lambda^N_a - \lambda^N_A ~.\label{eq23}
\end{equation}
The Fermi-levels $\lambda^Z_a and \lambda^N_a$ e.g. in the projectile
are defined by
\begin{equation}
     \lambda^N_a = \half \Bigl( M(Z_a,A_a+1)-M(Z_a,A_a-1)\Bigr) 
      -M(0,1) \label{eq20}
\end{equation}
and
\begin{equation}
 \lambda^Z_a = \half \Bigl( M(Z_a+1,A_a+1)-M(Z_a-1,A_a-1)\Bigr) 
      -M(1,1)~, \label{eq21}
\end{equation}
$M(Z,A)$ being the mass of the nucleus indicated. The model described in
\cite{W1} allows for Q-values where $Q_{gg}^P \not= -Q_{gg}^S$. Two
possible choices of this type are optional.
\par
The kind of information searched for the collective modes is governed
by the flag {\it inelas}, which has the value -1 if inelastic is to be
neglected, 0 if only giant resonances, 1 if only low-lying states, and
2 if both low- and high-lying states are to be included. If data for the
low-lying states are not found in the table one is asked to provide 
information to supplement the table. The values {\it inelas} equal to 3 
or 4, has the same significance as 1 and 2 respectively, but is used
in the {\it vms} version for automatic runs, where data are skipped 
if not found in the table. 

\subsection{Subroutine CAPTURE\_S}\label{CAPT}
In Fig.\ref{fig3} is shown the flowdiagram for the subroutine CAPTURE\_S, 
where the classical capture angular momentum is evaluated. The central 
part of this program is the integration routine from the NAG library and
the subroutine FUNCAP which evaluates the equations to be integrated. 
If so specified one may change the value of the damping,
change to a linear damping, or introduce a quartic term in the restoring 
force  for these modes as described in \cite{BDW}. The integration is 
made with {\it w} as variable, which measures time a unit that increases 
with distance to make a convenient treatment of Coulomb excitation. 
The experience is that only very strong dampings or non-linearities 
change the capture distance {\it rcapt} appreciably. They do however 
have a strong influence on the shapes of the two nuclei as they develop
after $t=w=0$. If additional output was asked for in the beginning of
the program, one may obtain tables which are convenient for plotting the 
shapes. The program keeps track of the overall energy and angular 
momentum conservation during integration. Because of neglected terms in 
the equations of motion no energy conservation can be expected after 
capture. 

\subsection{Subroutine PROBBAR\_S}\label{Prob}
In Fig.\ref{fig4} is shown the flowdiadram of the routine in which the 
probability distribution of barriers are computed.
The basis for this calculation is a trajectory with {\it r} 
equal to the capture distance {\it rcapt}. 
For the transfer reactions a Gaussian distribution
is assumed, with a width computed according to ref.\cite{W2}. The mean
energy shift from transfer as well as from the surface modes
is assumed to be taken into account by shifting the barrier to the 
capture distance. 


\par
For the inelastic modes each magnetic substate should be treated separately
if {\it phidot} is nonvanishing. If {\it phidot} is smaller than
{\it phidotmin} the m-substates are lumped together but given a width
{\it widspr} corresponding to the difference in "excitation energy" {\it alf} 
($= \hbar \omega -  m \dot \Phi$). Also, since a small value of {\it alf}
would give rize to high "excitation probabilities" {\it probb} and a nearly
Gaussian barrier distribution, states for which {\it alf} is smaller than
{\it esmin} are taken together with the transfer contributing to the 
total width (i.e. the grand total width reduced by those inelastic modes
which are treated separately). Also states with {\it probb} less than
{\it probmin} are neglected.
For the selected states (with higher excitation energy 
and sufficient excitation probability) a discrete Poisson distribution 
is assumed. 
\par
No explicit implementation is as yet made for deformed nuclei. They are 
treated as (soft) vibrators giving rise to a Gaussian barrier distribution,
which may be quite adequate.
\par
The result is stored in the {\it barprob} array for later use to 
calculate the capture probability in the TRANSM\_S subroutine.

\subsection{Subroutine TRANSM\_S}\label{trans}
This subroutine calculates in a straight forward way the capture 
probability by folding the barrier distribution {\it barprob} with
the WKB penetrability. The mass parameter in the WKB approximation 
is adjusted by a factor, $1/(1 + {\it delta})$ (cf. ref.\cite{W2} 
eq.(22)), to take into account the mass parameters of the 
collective modes. The prescription is not accurate and may be changed 
in the original input. The result for the capture probability at the 
classical capture should be close to 0.5.
\par
A problem arises for trajectories inside the capture distance. As was 
mentioned above the reactions close to (and inside) this distance are 
reduced with the philosophy that for the reflected trajectory the 
surface-surface distance has on the average been further apart than 
the grand average. For the trajectories that would surpass the nominal 
barrier there is, however, no provision to calculate the reactions
and the corresponding cross section is unaccounted for.
This point needs further consideration.

\subsection{Subroutine NZDISTRI\_S}\label{nzdis}
This subroutine calculates the distribution in $N$ and $Z$ as well
as the mean and spreads in energy and mean angular momentum for each 
partition due to transfer reactions.
The program evaluates the following quantities:
\begin{itemize}
\item $P_N = p0nn(ideln),~~~~ P_Z = p0zz(idelz), $
\item $\langle E^* \rangle _{\Delta N, \Delta Z} = eavdist(ideln,idelz), 
~~~~(\langle E^* \rangle _{0,0})^* = eavdist00,$
\item $(\sigma^2_{E^*})_{\Delta N,\Delta Z} = wid2dist(deln,delz),
~~~~(\sigma^2_{E_a^*})_{\Delta N,\Delta Z} = wid2adist(deln,delz),$
\item $\langle M \rangle _{\Delta N,\Delta Z} = avmdist(ideln,idelz),
~~~~(\langle M \rangle _{0,0})^* = avmdist00,$
\item $\langle M_a \rangle _{\Delta N,\Delta Z} = avmadist(ideln,idelz),
~~~~(\langle M_a \rangle _{0,0} = avmadist00.$
\end{itemize}
\par
In this connection it should be noted that ref.\cite{W2} contains 
a couple of misprints and errors. Thus Eq.(5.16) should read
\begin{equation}
      \langle E^*\rangle_{0,0}^* = \frac{P_{00}}{1-P_{00}}
      \langle E^*\rangle_{0,0}     ~,\label{eq3}
\end{equation}
while Eq.(5.19) should read
\begin{equation}
      \left(\sigma_{E^*}^2\right)_{0,0}^*
      =\frac{P_{00}}{1-P_{00}}\left(\left(\sigma_{E^*}^2\right)_{0,0}
      -\frac{p_0}{P_{00}-p_0}\langle E^*\rangle^2_{0,0}\right)
     ~.\label{eq4}
\end{equation}
\par
Finally Eq.(5.24) is incorrect. The quantity 
$\left(\sigma_{E^*_b}^2\right)_{\Delta N,\Delta Z}$ has to be evaluated
from Eq.(5.18) using the above results (\ref{eq1}) and (\ref{eq2}).

\section{Subroutine LEG1\_S}\label{leg1}
In jobtype 1 only average values of the reaction products for a given
trajectory are calculated. These are:
\begin{itemize}
\item Average excitation energy and spread of recoil, ejectile and total
kinetic energy ({\it eava, widea, eavaa, wideaa, tav}).
\item Average angular momentum of recoil and ejectile and total angular 
momentum loss ({\it avma, avmaa, avm}).
\item Average drift and spread of mass and charge ({\it drifta, wida, driftz,
widz}).
\item Average number of phonons in each of the collective modes ({\it 
ff0in(i,lam,n)}).
\end{itemize}
In addition one obtains the capture probability ({\it ppcapt}), the
elastic scattering angle ({\it thet}), the probability for elastic
scattering  and the absorptive potential ({\it ww}). The mean 
scattering angle for reactions and the corresponding differential 
cross sections are calculated by using an average trajectory between
the elastic trajectory and one which has suffered the mean mass and 
charge drift and the mean energy and angular momentum loss.
\par
A flowdiagram of LEG1\_S is given in Fig.\ref{fig5}. Since most of 
the quantities are evaluated in the main program the main
part of LEG1\_S is therefore instructions about the output. One has 
however the option to calculate the probability for kinetic energy 
loss larger than a given cutoff $E_0$ and with a finite energy 
resolution $\Delta$. This is evaluated in the subroutine TRESLIM\_S 
according to the expression
\begin{eqnarray}
      P(E^T>E_0) & = & \frac{p_0}{2\pi i}
      \int_{-\infty}^{\infty}\frac{\eto^{-i\beta^T E_0}}
      {\beta^T-i\delta_+}\exp\left\{-\half\Delta^2(\beta^T)^2\right\}
      \nonumber\\ & \times & 
      \left[\exp\left(\sum \eto^{-iQ_{gg}^{\tau}\beta^T}
      F^{(\tau)}(\beta^T)\right)\right] ~d\beta^T
       ~, \label{eq5}
\end{eqnarray}
which corresponds to Eq.(4.23) of ref.\cite{W1} for finite energy 
resolution. With finite energy resolution ( {\it tres} $\not= 0$ ) the 
reaction cross section contains some elastic cross section and the
total cross section diverges. If total cross sections are to be 
calculated one must therefore indicate the minimum (elastic) scattering 
angle ({\it scatmin}) to be included.

\subsection{Subroutine TRESLIM\_S}\label{treslim}
The subroutine TRESLIM\_S evaluates the expression (\ref{eq5}).

\section{Subroutine LEG2\_S}\label{leg2}
The jobtype 2 involves the calculation of the distribution in $N$ and $Z$
after the collision and related quantities. These are:
\begin{itemize}
\item Average excitation energy and the energy spread of each fragment 
(recoil or ejectile) and total kinetic energy loss for each partition.
\item Average angular momentum of each fragment (recoil or ejectile) 
and total angular momentum loss for each partition.
\end{itemize}
All these quantities are evaluated for transfer in NZDISTRI\_S 
(cf. Section \ref{nzdis}), and for inelastic scattering in the 
main program. A flowdiagram of LEG2\_S is given in Fig.\ref{fig6}.
\par
If only a single trajectory is asked for the above quantities are 
printed out together with the scattering angle and differential cross 
section for each partition. The latter quantities are evaluated in the
subroutine SCATANGLEAVX\_S on the basis of the calculated mean
energy and mean angular momentum loss for the partition in question.
\par
If total cross sections are asked for one may obtain printout and tables
of all these quantities for each of the trajectories contributing to the
total cross section. Since this may imply a very large output one is asked 
about which quantities and which partitions one wants in printout or in
tables.

\subsection{Subroutine SCATANGLEAVX\_S}\label{scatangleavx}
This subroutine evaluates the scattering angle for a reaction, where 
the mass, charge, energy and orbital angular momentum has been changed 
by given amounts. While for simple one-step processes the angle 
could be given by a simple average between the scattering angle in the 
entrance and exit channels, the scattering angle for a deep inelastic
reaction cannot be estimated this way. In the present program
the trajectory reaches the turning point in the entrance channel
and is gradually changed on the way out. The prescription for the
rate of the energy loss is taken from Fig.5 of ref.\cite{W2} to
take place over a range the order of unity of the parameter $w$ 
($w0-delw|w0+delw$) centered around $w = w0$.
The transfer of angular momentum and mass and charge is assumed 
to be faster (no polarization effect) and thus take place around
$w0/2$ with a width $delw/2$. The quantities $w0$ and $delw$ 
are input to the program. Standard values are $w0=1$ and $delw = .5$,
but the values are increased in the program (in a smooth way) for 
cases where the turning points in entrance and exit channels are 
very far apart.

\section{Subroutine LEG3\_S}\label{leg3}
In this jobtype energy distributions are evaluated summed over
all partitions and over all transferred angular momenta. One has three
options:\\
1. The distribution in total kinetic energy \\
2. The distribution in excitation energy of the ejectiles\\
3. The distribution in excitation energy of the recoil\\
For total cross sections one has the additional option to produce\\
4. A table of the double differential cross section as function of
angle and total kinetic energy loss to produce a Wilczynski Plot.
\par
In all cases a finite energy resolution {\it tres} is asked for, which 
must be greater than 1 MeV (FWHM). This facilitates the inclusion 
of the inelastic reactions. 
\par
With finite energy resolution the 
reaction cross section contains some elastic cross section and the
total cross section diverges. If total cross sections are to be 
calculated one should therefore indicate the minimum (elastic) scattering 
angle ({\it scatmin}) to be included.
\par
A flowdiagram of LEG3\_S is given in Fig.\ref{fig7}.
\subsection{Subroutine TDISTRI\_S}\label{tdistri}
The subroutine TDISTRI\_S evaluates the total kinetic energy distribution
according to
\begin{equation}
  P_{\Delta}(E^T)=  p_0\frac{1}{2\pi}\int_{-\infty}^{\infty}
      \eto^{-i\beta^T E^T}\exp\left\{-\half\Delta^2(\beta^T)^2\right\}
      z_{tr}(\beta^T)z_{in}^a(\beta^T)z_{in}^A(\beta^T) ~d\beta^T 
      ~.\label{eq6}
\end{equation}
with $\Delta = tres/2.35..$. The characteristic function due to transfer 
$Z_{tr}(\beta^T)=z_{tr}(\beta^T)/z_{tr}(0)$ is given in ref.\cite{W1} 
Eq.(4.27) i.e. 
\begin{eqnarray}
  z_{tr}(\beta^T) & = & \exp \left(
   \eto^{-iQ_{gg}^{\nu S}\beta^T}F^{\nu S}(\beta^T,\beta^T,0,0)
  +\eto^{iQ_{gg}^{\nu S}\beta^T}F^{\nu P}(\beta^T,\beta^T,0,0)\right.
     \label{eq7} \\   & ~ & ~~~~ + \left.
   \eto^{-iQ_{gg}^{\pi S}\beta^T}F^{\pi S}(\beta^T,\beta^T,0,0)
  +\eto^{iQ_{gg}^{\pi S}\beta^T}F^{\pi P}(\beta^T,\beta^T,0,0)\right)
     ~.\nonumber
\end{eqnarray}
For the excitation of the inelastic modes, $n$, the characteristic function
is given in ref.\cite{W1} Eqs.(2.54)-(2.57), where (2.57) should read 
\begin{equation}
      p_m^a(\omega_m,\mu_m)=\frac{2\pi a}{\hbar^2\ddot r_0}
      \left( \frac{\partial U^N_{aA}}{\partial r}\right)^2
      \alpha^2_m R_a^2 \left( Y_{\lambda_m \mu_m}
      ({\scriptstyle \frac{\pi}{2}},0)\right)^2
      \eto^{-(\omega_m-\mu_m\dot\Phi_0)^2 a/\ddot r_0}
     ~. \label{eq8}
\end{equation}
One thus finds the characteristic function 
$Z_{in}^a(\beta^T)=z_{in}^a(\beta^T)/z_{in}^a(0)$ with
\begin{equation}
  z_{in}^a(\beta^T)  = \exp\left(\sum_m p_m^a \eto^{(i\hbar \omega_m \beta^T
            -(\sigma_{gm} \beta^T)^2/2)}\right) ~,\label{eq9}
\end{equation}
and 
\begin{equation}
  p_m^a(\omega_m) = \sum_{\mu_m} p_m^a(\omega_m,\mu_m)~.\label{eq10}
\end{equation}
In Eq.\ref{eq9} is also included a width $\sigma_{gm}$ of the giant
resonances, leading to a mean square spread of the levels proportional to
the number of phonons in the mode.
\par
It is noted that the factor $p_0=p_{tr} p_a p_A$ in Eq. \ref{eq6} being 
the probability for elastic scattering arizes from the factors $1/z_{tr}(0)$,
$1/z_{in}^a(0)$ and $1/z_{in}^A(0)$.

\subsection{Subroutine EDISTA\_S}\label{edista}
This subroutine evaluates the energy distribution of ejectile or recoil
summed over all partitions. For the ejectile it is given by
\begin{equation}
  P_{\Delta}(E_a)=  p_{tr} p_a\frac{1}{2\pi}\int_{-\infty}^{\infty}
      \eto^{-i\beta_a E_a}\exp\left\{-\half\Delta^2(\beta_a)^2\right\}
      z_{tr}(\beta_a,0,0,0)z_{in}^a(\beta_a)~d\beta_a \label{eq11}
\end{equation}
with
\begin{eqnarray}
      z_{tr}(\beta_a,0,0,0) = &\exp\left( F^{\nu S}(\beta_a,0,0,0)
      +F^{\nu P}(\beta_a,0,0,0)\right.\nonumber \\   &~~~+ \left.
      F^{\pi S}(\beta_a,0,0,0)+F^{\pi P}(\beta_a,0,0,0)\right)
     ~,\label{eq12}
\end{eqnarray}
and similar for the recoil.

\subsection{Subroutine EDISTAPP\_S}\label{edistapp}
This subroutine evaluates the an approximate energy distribution of
ejectile or recoil i.e. an approximation to the result described in
the routine \ref{edista}. The approximation for the energy distribution
for transfer is based on the expression Eqs.(4.16)-(4.18) of ref. \cite{W1}.
One uses the the corresponding characteristic function:
\begin{equation}
     Z_{tr}(\beta_a) = \eto^{-(\beta_a \Delta)^2/2}
     \left( p_{tr}+(1-p_{tr})(1-i \beta_a \langle E \rangle _0 )^
      {(-\langle E_{tr} \rangle _0/\sigma_{0tr}^2)}\right) ~,\label{eq12a}
\end{equation}
where $\langle E_{tr} \rangle_0$ is the mean energy and $\sigma_{0tr}^2$ 
the mean square width for those transfer reactions (with probability 
$p_{tr}$) that do not lead to elastic scattering. The inelastic 
processes are included as in \ref{edista}.
\par
The approximation (\ref{eq12a}) is useful for the calculation of
evaporation probabilities in \ref{leg4}, where very many distributions 
are required. A comparison with the results of \ref{edista} gives an idea
about the accuracy that is achieved.

\subsection{Subroutine MAVTDIST\_S}\label{mavtdist}
In order to estimate the scattering angle for a given kinetic energy
loss, $E^T$, the corresponding average angular momentum loss should 
be evaluated. One finds from ref.\cite{W1} 
\begin{eqnarray}
  \langle M_b&+&M_B\rangle _{E^T} = \frac{p_0}{P_{\Delta}(E^T)}
      \frac{1}{2\pi}\int_{-\infty}^{\infty} ~d\beta^T \eto^{-i\beta^T E^T}
      \exp\left\{-\half\Delta^2(\beta^T)^2\right\}\nonumber \\
      &\times&\frac{\partial}{i\partial\zeta}\exp\left(\sum_{\tau} 
      \eto^{-iQ_{gg}^{\tau}\beta^T}F^{\tau}(\beta^T,\beta^T,\zeta,\zeta)
      +F^a(\beta^T,\zeta)+F^A(\beta^T,\zeta)\right)
      {\Big |}_{\zeta =0}\nonumber \\      
      & &~~~~~~~~~~~= \frac{p_0}{P_{\Delta}(E^T)}
      \frac{1}{2\pi}\int_{-\infty}^{\infty} ~d\beta^T 
      \eto^{-i\beta^T E^T}\exp\left\{-\half\Delta^2(\beta^T)^2\right\}
      \label{eq13} \\
      &\times&\exp\left(\sum_{\tau}\eto^{-iQ_{gg}^{\tau}\beta^T}
      F^{\tau}(\beta^T,\beta^T,0,0)
      +F^a(\beta^T,0)+F^A(\beta^T,0)\right) \nonumber \\
      &\times&\frac{\partial}{i\partial\zeta}
      \left(\sum_{\tau}\eto^{-iQ_{gg}^{\tau}\beta^T}
      F^{\tau}(\beta^T,\beta^T,\zeta,\zeta)+F^a(\beta^T,\zeta)
      +F^A(\beta^T,\zeta)\right){\Big |}_{\zeta =0}\nonumber
\end{eqnarray}
with
\begin{eqnarray}
      \frac{\partial F^{\tau}}{i\partial\zeta}{\Big |}_{\zeta =0}
       & = &  k_\nu(r_0)\eto^{L_\nu^{(0)}-(J_{\nu}^{(0)})^2} 
      \left[ 4 \frac{\dot \Phi_0 c_3^2(b_a^{\nu}+b_A^{\nu})}
      {B_{\nu}\omega_{0\nu}}
      {\rm G}_2(J_{\nu}^{(0)}-i\hbar\omega_{0\nu}B_{\nu}\beta^T/2)\right.
      \label{eq14}\\
       & + &\left.\left( \frac{(b_A^{\nu}-b_a^{\nu})}{D_{\nu}}
      \frac{k_{\bot}}{\kappa_{\nu}}+ 
       2 J_{\nu}^{(0)}\frac{\dot \Phi_0 c_3^2(b_a^{\nu}+b_A^{\nu})}
      {B_{\nu}\omega_{0\nu}}\right)
      {\rm G}_1(J_{\nu}^{(0)}-i\hbar\omega_{0\nu}B_{\nu}\beta^T/2)\right]
      \nonumber 
\end{eqnarray}
and
\begin{equation}
      \frac{\partial F^{a}}{i\partial\zeta}{\Big |}_{\zeta =0}
      = \sum_m q_m^a \eto^{(i\hbar \omega_m \beta^T
            -(\sigma_{gm} \beta^T)^2/2)} \label{eq15}
\end{equation}
with
\begin{equation}
   q_m^a(\omega_m) = \sum_{\mu_m} \mu_m p_m^a(\omega_m,\mu_m)~.\label{eq16}
\end{equation}

\subsection{Subroutine DRIFTNZTDIS\_S}\label{driftnztdis}
In order to estimate the scattering angle one also needs the average 
drift of neutrons and protons for given kinetic energy loss. 
One finds from ref.\cite{W1} the average neutron drift
\begin{eqnarray}
     \langle \Delta N\rangle _{E^T} &=& \frac{p_0}{P_{\Delta}(E^T)}
      \frac{1}{2\pi}\int_{-\infty}^{\infty} ~d\beta^T 
      \eto^{-i\beta^T E^T}\exp\left\{-\half\Delta^2(\beta^T)^2\right\}
      \label{eq17} \\
      &\times&\exp\left(\sum_{\tau}\eto^{-iQ_{gg}^{\tau}\beta^T}
      F^{\tau}(\beta^T,\beta^T,0,0)
      +F^a(\beta^T,0)+F^A(\beta^T,0)\right) \nonumber \\
      &\times&\left(\eto^{-iQ_{gg}^{\nu P}\beta^T}   
       F^{\nu P}(\beta^T,\beta^T,0,0)-\eto^{-iQ_{gg}^{\nu S}\beta^T}
      F^{\nu S}(\beta^T,\beta^T,0,0)
      \right)\nonumber
\end{eqnarray}
and similar for the proton drift.

\subsection{Subroutine WILSH\_S}\label{wilsh}
The scattering angle for given total kinetic energy loss $E^T$ is evaluated
from the mean angular momentum loss and mass and charge drifts for this
energy in the routine \ref{scatangleavx}. One may obtain an approximate 
differential cross section corresponding to each trajectory by using
the elastic geometric cross section as is done under option 1. 
In order to obtain the true angular distributions one must consider
the actual correlation between scattering angle and impact parameter and
sum over several trajectories. In the present subroutine this is done
by interpolation, producing a table of the cross section $d\sigma$ 
associated with the turning point between $r$ and $r+dr$ i.e.
\begin{equation}
   \frac{d\sigma}{dr} = \rho (r) \frac{d\rho}{dr} P(E^T) P_{capt} \nonumber
\end{equation}
and a corresponding dense table of the scattering angle $\theta'(r)$.
To smooth out the rainbow angles the cross section $\rho d\rho 
/({\rm sin}\theta d\theta )$ is folded with a Gaussian distribution
$\exp ((\theta - \theta'(r))/\Delta )^2/2)$ of width $\Delta$, which is
an input to the routine. The resulting expression for 
the differential cross section is
\begin{equation}
   \frac{d\sigma}{d\Omega} = \frac{1}{{\rm sin}\theta}
      \frac{1}{{\sqrt{2\pi}\Delta}}
      \int _{-\infty}^{\pi}~dr \frac{d\sigma}{dr}
      \exp\frac{(\theta - \theta'(r))^2}{2(\Delta )^2}~.\label{eq18}
\end{equation}
A flow diagram of the routine is given in Fig.\ref{fig8}.
\subsection{Subroutine WILSH\_PLOT\_S}\label{wilsh_plot}
This subroutine produces a command file to make a Wilczynski plot
on the basis of the table produced in WILSH\_S. To run this command
file one needs the program NANNI.GRAPH.

\section{Subroutine LEG4\_S}\label{leg4}
In this jobtype energy distributions are evaluated for each partition
summed over all transferred angular momenta. In calculating energy
distributions for partitions far from the entrance channel one meets the
problem that the nuclei here have less binding energy than one would
assume from the simple model of multinucleon transfer between two 
fixed mean fields. In the present program it is assumed that the
transfer of nucleons is associated with a change in total kinetic
energy which follows the simple model with a constant minimum 
Q-value per step. After the collision the transferred nucleons find
themselves in a nucleus which most often offers less binding than was
assumed. This is not very serious since the mean excitation energy 
associated with multinucleon transfer is usually quite high. The 
excitation energy of the final nuclei is simply assumed to be less by 
the difference between the model binding energies and the actual binding 
energies of the recoil and the ejectile in question. For the part of the
distribution where there are no bound states available, the transfer is 
assumed to be blocked and the transfer probabilities are modified 
accordingly. All these calculations of multinucleon transfer are made with
approximate energy distributions based on gamma-distributions with the
correct mean excitation energies and spreads.
\par
For reactions with a small net number of transfered nucleons, a more 
accurate energy distribution should be obtainable. It is numerically 
calculated in option 3 for selected partitions, together with the mean 
angular momentum for a given excitation energy (and for this partition).
\par
Jobtype 4 offers 5 options:\\
1. Approximate distribution in excitation energy of all recoils or 
ejectiles corrected for the actual binding energies as described above \\
2. Approximate distribution in excitation energy of recoils or ejectiles
after correction for actual masses and after evaporation\\
3. The detailed distribution in excitation energy of selected recoils
or ejectiles\\
4. The detailed distribution in total kinetic energy of selected 
partitions\\
For total cross sections one has the additional option to produce\\
5. A table of the double differential cross section as function of
mass or charge and total kinetic energy loss to produce a plot.
\par
A flowdiagram of LEG4\_S is given in Fig.\ref{fig9}. This routine
evaluates the distribution over N and Z and the corresponding mean
excitation energies for each partition to be used in the various options,
which are, however, evaluated in separate subroutines described below.   
\subsection{Subroutine LEG41\_S}\label{leg41}
In this routine the actual binding energies (masses) of the product nuclei
are read in the mass table and compared to the masses assumed in the model.
A flowdiagram of LEG41\_S is given in Fig.\ref{fig10}.
In the standard option for the ground state to ground state Q-values
Eqs.(\ref{eq22})-(\ref{eq23}) the model mass of an ejectile is 
\begin{eqnarray}
   M_{app}(Z_a+\Delta Z,A_a+\Delta N)= M_{app}(Z_a,A_a)
      &+& (M(1,1)-\lambda^Z_a) \Delta Z \\ \nonumber
      &+& (M(0,1)-\lambda^N_a) \Delta N ~,\label{eq19}
\end{eqnarray}
where the mass $M_{app}$ is given by
\begin{eqnarray}
   M_{app}(Z_a,A_a) &=& {\mbox{$\frac{1}{4}$}}\Bigl( M(Z_a+1,A_a+1)
      +M(Z_a-1,A_a-1)\\ \nonumber
       &+&  M(Z_a,A_a+1)+M(Z_a,A_a-1)\Bigr) ~.\label{eq19a}
\end{eqnarray}
For other Q-value options the model mass is changed correspondingly.
\par
The distribution in excitation energy of the fragments is estimated
by a gamma-distribution with the calculated mean energy and spread
for the partition in question. This distribution over N and Z is 
corrected by the part of the distribution that lies below the actual
ground state mass in the subroutine ENZCOR\_S. 

\subsection{Subroutine LEG42\_S}\label{leg42}
In this routine the mass and charge distributions after evaporation
is estimated. It is made on the basis of the estimated distribution 
in excitation energy of the fragments as given by a gamma-distribution 
with the calculated mean energy and spread for the partition in question. 
The distribution over N and Z is then corrected in the routine ENZCOR\_S
for the actual masses as described in Section \ref{leg42}. Finally
the evaporation of neutrons is estimated in the subroutine EVAP\_S.
A flowdiagram of LEG42\_S is given in Fig.\ref{fig11}.
\subsection{Subroutine LEG43\_S}\label{leg43}
In this routine the distribution in excitation energy of selected
ejectiles or recoils are calculated numerically. For an ejectile
of charge number $Z_a+\Delta Z$ and mass number $A_a+\Delta N+\Delta Z$
one finds the distribution
\begin{eqnarray}
 & &P_{\Delta N \Delta Z}(E^*_a) = \frac{p_{tr} p_a}{P_{\Delta N \Delta Z}}
      \frac{1}{2\pi}\int_{-\infty}^{\infty}~d\beta_a
      \eto^{-i\beta_a E_a}\label{eq25} \\
      &\times &\exp\Bigl( \sum_{\tau}F^{\tau}(\beta_a,0,0,0)
      +F^{a}(\beta_a,0)\Bigr) B_{\Delta N}(F^{\nu P},F^{\nu S}) 
      B_{\Delta Z}(F^{\pi P},F^{\pi S}) ~.\nonumber 
\end{eqnarray}
In this expression $F^{\tau}=F^{\tau}(\beta_a,0,0,0)$ 
is given by ref.\cite{W1} Eq.(2.38) and $F^a=F^a(\beta_a,0)$ is 
given by Eq.(2.55) with $p^a_m$ as defined in (\ref{eq8}) above.
The results are compared to the results of the approximation,
where a gamma-distribution with the correct mean energy and spread 
is used. This comparison is a help in estimating the accuracy of 
the approximate treatments in LEG41\_S and LEG42\_S.

A flowdiagram of LEG43\_S is given in Fig.\ref{fig12}.
\par
It should be noted that because the distribution in excitation energy
is not corrected for the actual mass of the fragment as
described in Section \ref{leg41} above, the results can
only be trusted for small values of $\Delta N$ and $\Delta Z$.

\subsection{Subroutine LEG44\_S}\label{leg44}
In this routine the distribution in total kinetic energy for selected
partitions are calculated numerically. For a partition, where the ejectile
has charge number $Z_a+\Delta Z$ and mass number $A_a+\Delta N+\Delta Z$
one first calculates the distribution in total excitation energy
\begin{eqnarray}
 & &P_{\Delta N \Delta Z}(E^*) = \frac{p_{0}}{P_{\Delta N \Delta Z}}
      \frac{1}{2\pi}\int_{-\infty}^{\infty}~d\beta
      \eto^{-i\beta E^*}\label{eq26} \\
      &\times &\exp\Bigl( \sum_{\tau}F^{\tau}(\beta ,\beta ,0,0)
      +F^{a}(\beta_a,0)+F^{A}(\beta_a,0)\Bigr) 
      B_{\Delta N}(F^{\nu P},F^{\nu S}) 
      B_{\Delta Z}(F^{\pi P},F^{\pi S}) ~,\nonumber 
\end{eqnarray}
where $F^{\tau}=F^{\tau}(\beta ,\beta ,0,0)$ is given by ref.\cite{W1} 
Eq.(2.38). The distribution in the total kinetic energy, $T$, is obtained
from here since 
\begin{equation}
      T=E(\alpha )-E^*-\Delta N Q_{gg}^{\nu S}-
      \Delta Z Q_{gg}^{\pi S}~.\label{eq27}
\end{equation}
This distribution is not expected to be very sensitive to the problem
that the actual masses of the reaction products differs from the 
model-masses (\ref{eq19}), and might therefore be trusted in a much
larger $\Delta N,\Delta Z$ range.
A flowdiagram of LEG44\_S is given in Fig.\ref{fig13}.
\par
In order to obtain a realistic angular distribution of the fragments
one should calculate the mean loss of orbital angular momentum for the
partition in question $\Delta L = -\langle M_a + M_A \rangle $. 
This is done in the subroutine MAVNZDIS\_S.
On the basis of $\Delta L, T, \Delta N$ and $\Delta Z$ the scattering 
angle is calculated in the subroutine SCATANGLEAVX\_S described in 
Section \ref{scatangleavx}.

\subsection{Subroutine MAVNZDIS\_S}\label{mavnzdis}
In order to estimate the scattering angle for a given partition 
$(\Delta N,\Delta Z)$ and a given kinetic energy loss, the average 
angular momentum loss for given total excitation energy, $E^*$, should 
be evaluated. One finds from ref.\cite{W1} 
\begin{eqnarray}
\lefteqn{\langle M_b + M_B \rangle _{\Delta N\Delta Z,E^*} } \nonumber \\
 \lefteqn{ =     \frac{p_0}{P_{\Delta N\Delta Z}(E^*)}
      \frac{1}{2\pi}  \int_{-\infty}^{\infty} ~d\beta~\eto^{-i\beta E^*}
      \frac{\partial}{i\partial\zeta} \biggl( 
      B_{\Delta N}(F^{\nu P},F^{\nu S}) B_{\Delta Z}(F^{\pi P},F^{\pi S}) 
      }\nonumber \\
&\times  \exp \Bigl( \sum_{\tau}F^{\tau}(\beta ,\beta ,\zeta ,\zeta )
      +F^a(\beta ,\zeta ) +F^A(\beta ,\zeta ) \Bigr) \biggr) 
      {\Big |}_{\zeta =0}  \nonumber \\
      \lefteqn{ =  \frac{p_0}{P_{\Delta N\Delta Z}(E^*)}
      \frac{1}{2\pi}\int_{-\infty}^{\infty} ~d\beta~ 
      \eto^{-i\beta E^*}     
     \exp \Bigl( \sum_{\tau} F^{\tau}(\beta ,\beta,0,0)
      +F^a(\beta,0)+F^A(\beta,0)\Bigr) }\nonumber \\
     & \left( B_{\Delta N+1}(F^{\nu P},F^{\nu S})
      B_{\Delta Z}(F^{\pi P},F^{\pi S})
      \frac{\partial F^{\nu S}}{i\partial \zeta}
       + B_{\Delta N-1}(F^{\nu P},F^{\nu S})
      B_{\Delta Z}(F^{\pi P},F^{\pi S})
      \frac{\partial F^{\nu P}}{i\partial\zeta}  \right.   \nonumber \\
      &+ B_{\Delta N}(F^{\nu P},F^{\nu S})
      B_{\Delta Z +1}(F^{\pi P},F^{\pi S})
      \frac{\partial F^{\pi S}}{i\partial\zeta}
       + B_{\Delta N}(F^{\nu P},F^{\nu S})
      B_{\Delta Z -1}(F^{\pi P},F^{\pi S})
      \frac{\partial F^{\pi P}}{i\partial\zeta}  \nonumber \\
      &\left. \frac{\partial F^a(\beta,\zeta)}{i\partial\zeta}+
      \frac{\partial F^A(\beta,\zeta)}
      {i\partial\zeta} \right) {\Big |}_{\zeta =0}~,    \label{eq28} 
\end{eqnarray}
where the partial derivatives of $F^{\tau}$ are given in Eqs.(\ref{eq14})
and (\ref{eq15}) and $B_n(x,y)$ are the Bessel distributions given in
ref.\cite{W1} Eq.(5.2).
\subsection{Subroutine LEG45\_S}\label{leg45}
This routine produces tables for contour plotting of the cross section
as a function of the total kinetic energy and of the charge or the mass
of the ejectile. The energy distribution is calculated according to the
approximation in subroutine PROBDTAB\_S in terms of the gamma distribution 
for each mass and charge partition.

\begin{thebibliography}{19}
\bibitem{W1}          A.Winther
                      {\em Nucl.\ Phys.} {\bf A572} (1994) 191
\bibitem{W2}          A.Winther
                      {\em Nucl.\ Phys.} {\bf A594} (1995) 203
\bibitem{AW-81}       R.\"{O}. Aky\"{u}z and A. Winther,
                      in {\em Proceedings of the Enrico Fermi
                      International School of Physics, 1979,}\/
                      "Nuclear Structure and Heavy-Ion Reactions",
                      Eds.\ R.A. Broglia, C.H. Dasso, and R. Ricci
                      (Amsterdam: North-Holland Publ.~Co., 1981) 491
\bibitem{BDW}         R.A. Broglia, C.H. Dasso and A. Winther 
                      in Proceedings of the Enrico Fermi International 
                      School of Physics 1979 (ed. R. A. Broglia, 
                      C. H. Dasso and R. Ricci) North Holland Publ. Co., 
                      Amsterdam 1981
\bibitem{BW}          R.A. Broglia and A. Winther: Heavy Ion Reactions 
                      Parts I \& II, Frontiers in Physics Vol. 84, 
                      Addison-Wesley Publ. Co., 1991

\end{thebibliography}
\begin{figure}[p]
\caption{Flowdiagram for the main program. The numbers indicated in 
 circles refer to statement numbers in the program. The various boxes 
 indicate input(prompts), output, calculations, and branches. The 
 quantities in ellipses indicate calls to subroutines.\label{fig1}}
\end{figure}
\begin{figure}[p]    
\caption{Flowdiagram for the subroutine QVALETC\_S. The numbers indicated in 
 circles refer to statement numbers in the program. The various boxes 
 indicate input(prompts), output, calculations, and branches. The 
 quantities in ellipses indicate calls to subroutines.\label{fig2}}
\end{figure}
\begin{figure}
\caption{Flowdiagram for the subroutine CAPTURE\_S. The numbers indicated in 
 circles refer to statement numbers in the program. The various boxes 
 indicate input(prompts), output, calculations, and branches. The 
 quantities in ellipses indicate calls to subroutines.\label{fig3}}
\end{figure}
\begin{figure}[p]
\caption{Flowdiagram for the subroutine PROBBAR\_S. The numbers indicated in 
 circles refer to statement numbers in the program. The various boxes 
 indicate input(prompts), output, calculations, and branches. The 
 quantities in ellipses indicate calls to subroutines.\label{fig4}}
\end{figure}
\begin{figure}[p]
\caption{Flowdiagram for the subroutine LEG1\_S. The numbers indicated in 
 circles refer to statement numbers in the program. The various boxes 
 indicate input(prompts), output, calculations, and branches. The 
 quantities in ellipses indicate calls to subroutines.\label{fig5}}
\end{figure}
\begin{figure}[p]
\caption{Flowdiagram for the subroutine LEG2\_S. The numbers indicated in 
 circles refer to statement numbers in the program. The various boxes 
 indicate input(prompts), output, calculations, and branches. The 
 quantities in ellipses indicate calls to subroutines.\label{fig6}}
\end{figure}
\begin{figure}[p]
\caption{Flowdiagram for the subroutine LEG3\_S. The numbers indicated in 
 circles refer to statement numbers in the program. The various boxes 
 indicate input(prompts), output, calculations, and branches. The 
 quantities in ellipses indicate calls to subroutines.\label{fig7}}
\end{figure}
\begin{figure}[p]
\caption{Flowdiagram for the subroutine WILSH\_S. The numbers indicated in 
 circles refer to statement numbers in the program. The various boxes 
 indicate input(prompts), output, calculations, and branches. The 
 quantities in ellipses indicate calls to subroutines.\label{fig8}}
\end{figure}
\begin{figure}[p]
\caption{Flowdiagram for the subroutine LEG4\_S. The numbers indicated in 
 circles refer to statement numbers in the program. The various boxes 
 indicate input(prompts), output, calculations, and branches. The 
 quantities in ellipses indicate calls to subroutines.\label{fig9}}
\end{figure}
\begin{figure}[p]
\caption{Flowdiagram for the subroutine LEG41\_S. The numbers indicated in 
 circles refer to statement numbers in the program. The various boxes 
 indicate input(prompts), output, calculations, and branches. The 
 quantities in ellipses indicate calls to subroutines.\label{fig10}}
\end{figure}
\begin{figure}[p]
\caption{Flowdiagram for the subroutine LEG42\_S. The numbers indicated in 
 circles refer to statement numbers in the program. The various boxes 
 indicate input(prompts), output, calculations, and branches. The 
 quantities in ellipses indicate calls to subroutines.\label{fig11}}
\end{figure}
\begin{figure}[p]
\caption{Flowdiagram for the subroutine LEG43\_S. The numbers indicated in 
 circles refer to statement numbers in the program. The various boxes 
 indicate input(prompts), output, calculations, and branches. The 
 quantities in ellipses indicate calls to subroutines.\label{fig12}}
\end{figure}
\begin{figure}[p]
\caption{Flowdiagram for the subroutine LEG44\_S. The numbers indicated in 
 circles refer to statement numbers in the program. The various boxes 
 indicate input(prompts), output, calculations, and branches. The 
 quantities in ellipses indicate calls to subroutines.\label{fig13}}
\end{figure}

\end{document}
